%\documentclass{beamer}
%\usepacakge[UTF8]{ctex}
%----------------------------------------------------------------------------------------
%    PACKAGES AND THEMES
%----------------------------------------------------------------------------------------

\documentclass[aspectratio=169,xcolor=dvipsnames]{beamer}
\usepackage[UTF8]{ctex}
\usetheme{SimplePlus}
\setbeamertemplate{footline}{}

\usepackage{hyperref}
\usepackage{graphicx} % Allows including images
\usepackage{booktabs} % Allows the use of \toprule, \midrule and \bottomrule in tables

\usepackage{mathpartir}

\usepackage{tikz-cd}

\usepackage{fontspec}

\newcommand\bbD{\mathbb{D}}
\newcommand\bbI{\mathbb{I}}
\newcommand\bbF{\mathbb{F}}
\newcommand\JEq{\mathtt{JEq}}
\newcommand\FTy{\mathtt{FTy}}
\newcommand\comp{\mathtt{comp}}
\newcommand\fst{\mathtt{fst}}
\newcommand\snd{\mathtt{snd}}
\newcommand\Ty{\mathtt{Ty}}
\newcommand\Ter{\mathtt{Ter}}
\newcommand\dM{\mathsf{dM}}
\newcommand\ttp{\mathtt{p}}
\newcommand\ttq{\mathtt{q}}
\newcommand\bfone{\mathbf{1}}
\newcommand\Set{\mathbf{Set}}
\newcommand\DM{\mathbf{DM}}
\newcommand\Id{\mathtt{Id}}
\newcommand\refl{\mathtt{refl}}
\newcommand\C{\mathbf{C}}
\newcommand\psC{\widehat{\mathbf{C}}}
\newcommand\extop{{\scalebox{0.4}[1]{\_}}.{\scalebox{0.8}[1]{\_}}}
\newcommand{\shortminus}{\raisebox{0.3ex}{\scalebox{0.8}[1.0]{-}}}
\newcommand{\mathshortminus}{\mathbin{\text{\normalfont\shortminus}}}
\newcommand{\ul}{{\scalebox{0.4}[1.0]{\_}}}

\DeclareMathOperator\FreeDeMorgan{FreeDeMorgan}
\DeclareMathOperator\obj{obj}
\DeclareMathOperator\y{y}

\usepackage{unicode-math}
%\renewcommand{\square}{□}

%----------------------------------------------------------------------------------------
%    TITLE PAGE
%----------------------------------------------------------------------------------------

\title{区间与面格}
%\subtitle{Subtitle}

%\author{Pin-Yen Huang}

%\institute
%{
%    Department of Computer Science and Information Engineering \\
%    National Taiwan University % Your institution for the title page
%}
%\date{\today} % Date, can be changed to a custom date
\date{}

%----------------------------------------------------------------------------------------
%    PRESENTATION SLIDES
%----------------------------------------------------------------------------------------

\begin{document}

\begin{frame}
    % Print the title page as the first slide
    \titlepage
\end{frame}

\begin{frame}{背景}
\begin{itemize}
  \item 任何立方集$X$都可以看作任何上下文$\Gamma$中的一个类型. 我们可以取$X'\in\Ty(\Gamma)=\obj(\Set^{(\int\Gamma)^{op}})$ 且 $X'(I,\rho)\triangleq X(I)$
  \item 我们可以把项$\Gamma\vdash i:\bbI$和$\Gamma\vdash\varphi:\bbF$解释为$\Ter(\Gamma\vdash\bbI)$和$\Ter(\Gamma\vdash\bbF)$的元素
  \item 上下文扩展 $\Gamma.\bbI$ 就是 CwF 中的上下文扩展操作
\end{itemize}
(注: 下文把 $\bbI$ 和 $\bbF$ 解释为立方集, 所以把它们当作类型使用时, 严格来说应该使用$\bbI'$和$\bbF'$. $\Ter(\Gamma\vdash\bbI')$, $\Ter(\Gamma\vdash\bbF')$, $\Gamma.\bbI'$)
\end{frame}

\begin{frame}[fragile,shrink]{$\bbI$}
  \begin{itemize}
    \item 把$\bbI$解释为立方集: $\bbI(I)\triangleq\dM(I)$, 态射映为对应的德摩根代数同态
    \item 对$\bbI$的另一种刻画是单元素集对应的 Yoneda 嵌入 $\y\{i\}=\square(\ul,\{i\})$
  \end{itemize}
  对每一个 $I$, 得到一个德摩根代数$\bbI(I)$, 把上面的结构统一记作$0,1,\ul\lor\ul,\ul\land\ul,\neg\ul$, 根据情况确定到底对应哪一个$I$. 有了这些数据, 我们可以定义自然变换 $\lor:\bbI\times\bbI\to\bbI$\[ \lor_{I}(r,s) \triangleq r \lor s \]
  对$f\in\square(J,I)=I\to\dM(J)$, 记$\bar{f}\triangleq\bbI(f):\dM(I)\to\dM(J)$
  \[
  \begin{tikzcd}
  \bbI(I)\times\bbI(I) \arrow[d,"{\lor_{I}}"'] \arrow[r,"{(\bar{f},\bar{f})}"] & \bbI(J)\times\bbI(J) \arrow[d,"{\lor_{J}}"]\\
  \bbI(I) \arrow[r,"\bar{f}"'] & \bbI(J)
  \end{tikzcd}
  \]
  由$\bar{f}(r\lor s) = \bar{f}(r)\lor\bar{f}(s)$, 以上图表交换, 我们确实定义了一个自然变换.
\end{frame}

\begin{frame}[shrink]{$\bbF$}
  立方类型论中的$\bbF$是$\{(i=0),(i=1)\mid i\in\bbD\}$上的自由分配格, 并且模去关系$(i=0)\land(i=1)=0_{\bbF}$. 它的元素(也称为面公式)可以记为
  \[
  \varphi,\psi ::= 0_{\bbF} \mid 1_{\bbF} \mid i=0 \mid i=1 \mid \varphi\land\psi \mid \varphi\lor\psi \quad(i\in\bbD)
  \]
  我们把$\bbF$解释为立方集($I\in\obj\square$, $f\in\square(J,I)=I\to\dM(J)$)
  \begin{itemize}
    \item $\bbF(I)$ 是仅涉及$I$中名字的面公式集合
    \item $\bbF(f):\bbF(I)\to\bbF(J)$, 对于$\bbF(I)$中的面公式, 它把名字 $i\in I$ 替换为 $f(i)\in\dM(J)$, 并进行以下项重写, 这样我们就得到$\bbF(J)$中的面公式
%\[
%\begin{aligned}
%(r \lor s) = 1 \quad & \mapsto \quad (r = 1) \lor (s = 1) \\
%(r \lor s) = 0 \quad & \mapsto \quad (r = 0) \land (s = 0) \\
%(r \land s) = 1 \quad & \mapsto \quad (r = 1) \land (s = 1) \\ 
%(r \land s) = 0 \quad & \mapsto \quad (r = 0) \lor (s = 0) \\
%(1 - r) = 1 \quad & \mapsto \quad r = 0 \\
%(1 - r) = 0 \quad & \mapsto \quad r = 1
%\end{aligned}
%\]
\[
\begin{array}{rclcrcl}
(r \lor s) = 1 & \mapsto & (r=1) \lor (s=1) & \quad & (r \lor s) = 0 & \mapsto & (r=0) \land (s=0) \\
(r \land s) = 1 & \mapsto & (r=1) \land (s=1) & \quad & (r \land s) = 0 & \mapsto & (r=0) \lor (s=0) \\
(1-r) = 1 & \mapsto & r=0 & \quad & (1-r) = 0 & \mapsto & r=1
\end{array}
\]
  \end{itemize}
  对$\varphi\in\Ter(\Gamma\vdash\bbF)$, 我们定义限制上下文$\Gamma,\varphi$为如下立方集
  \[
  (\Gamma,\varphi)(I)\triangleq\{\rho\in\Gamma(I)\mid\varphi(I,\rho)=1_{\bbF}\}
  \]
\end{frame}

\begin{frame}[shrink]{解释}
  我们在立方集范畴 $\widehat{□}\triangleq\Set^{\square^{op}}$ 中解释立方类型论. 我们在前面展示过怎样从预层范畴得到 CwF. 我们定义 $\Ty(\Gamma)\triangleq\obj\Set^{{(\int\Gamma)}^{op}}$, 前面我们说过, 任意立方集都可以看作 $\Gamma$ 中的一个类型, $\bbF$ 也一样, 我们有 $\bbF'\in\obj\Set^{{(\int\Gamma)}^{op}}$. 对$\varphi\in\Ter(\Gamma\vdash\bbF')\triangleq\hom(1,\bbF')$, $\varphi$ 是一个自然变换, 对$(I,\rho)\in\int\Gamma$, $I\in\obj\square$, $\rho\in\Gamma(I)$, $\varphi_{(I,\rho)}:1(I,\rho)\to\bbF'(I,\rho)$, 由$\bbF'$的定义, $\varphi(I,\rho)\triangleq\varphi_{(I,\rho)}(\ast) \in \bbF'(I,\rho)=\bbF(I)$.
\end{frame}

\end{document}